\section{Discussion and extensions}

\paragraph{Semi-parametric approaches.}
\label{sec:semi-param}
\label{sec:master-results-sem}

Our analysis separates likelihood-based approaches using all the labels, as
in~\eqref{eq:maximum-likelihood}, and majority vote estimators, which are
robust but have slower convergence, as in
Proposition~\ref{prop:lowd-misspcfd-majority-vote-m}.
%
While we mainly focus on efficiency bounds for 
procedures we view as closer to practice,
%---fitting a model based on available data---
in this section, we
present a few extensions to semiparametric approaches, more
as an amuse-bouche to motivate further efficiency investigations
rather than as a final word.
%% We thus develop a few convergence results into which we can substitute
%% semiparametric estimators.
%
In distinction from standard results in semiparametric theory
(e.g.~\cite[Ch.~25]{VanDerVaart98} or~\cite{BickelKlRiWe98}), the results
require little more than consistent estimation of the links $\sigma_j\opt$
to recover (rate optimal) $1/m$ scaling in the asymptotic covariance.

Let $\vec{\sigma} = (\sigma_1, \ldots, \sigma_m)$
be potentially distinct labeler link functions
(with $\sigma_j\opt$ denoting the true links).
Then for
$(x, y) \in \R^d \times \{\pm 1\}^m$ define the link-based loss
\begin{equation*}
  \loss_{\vec{\sigma}, \theta}(y \mid x)
  \defeq \frac{1}{m} \sum_{j = 1}^m \loss_{\sigma_j,\theta}(y_j \mid x).
\end{equation*}
%
To eliminate ambiguity in $\vec{\sigma}$,
we normalize $\ltwo{\theta\opt} = 1$,
and then define the population loss
\begin{equation*}
	L(\theta, \vec{\sigma})
	\defeq \frac{1}{m} \sum_{j = 1}^m
	\E[\loss_{\sigma_j, \theta}(Y_j \mid X)].
\end{equation*}
For the sequence $\{\vec{\sigma}_n\} \subset \fsymlinksp$
of (estimated) links and data $(X_i,
Y_i)$, $Y_i = (Y_{i1}, \ldots, Y_{im})$, we define the semi-parametric
estimator
\begin{align*}
	\spe \defeq \argmin_{\theta \in \R^d}
	\frac{1}{n} \sum_{i = 1}^n
	\loss_{\vec{\sigma}_n, \theta}(Y_i \mid X_i).
\end{align*}

Both consistency and
asymptotic normality hold under appropriate convergence of
the link functions (see Appendix~\ref{sec:semi-supp} for details).
Define the distance $d_{\flink}$ on $\R^d
\times \flink$ by
\begin{align}
  d_{\flink}\left((\theta_1, \vec{\sigma}_1), (\theta_2, \vec{\sigma}_2)
  \right)
  \defeq
  \ltwo{\theta_1 - \theta_2}
  + \LtwoPbold{\vec{\sigma}_1(-Y\<X, u\opt\>)
    - \vec{\sigma}_2(-Y\<X, u\opt\>)}.
  \label{eqn:link-distance}
\end{align}
%We make the following assumption.
%\begin{assumption}
%	\label{assumption:sp-general}
%	The links $\vec{\sigma}\opt \in \fsymlinksp$ are normalized so
%	that $\P(Y_j = y \mid X = x) = \sigma\opt_j(y \<x, u\opt\>)$, and
%	the sequence $\{\vec{\sigma}_n\} \subset \fsymlinksp$ is consistent:
%	\begin{equation*}
%		\LtwoPbold{\vec{\sigma}_n(-Y\<X, u\opt\>)
%			- \vec{\sigma}\opt(-Y\<X, u\opt\>)}
%		\cp 0.
%	\end{equation*}
%	Additionally, the mapping $(\theta, \vec{\sigma}) \mapsto
%	\nabla_\theta^2 \poploss(\theta,
%	\vec{\sigma})$ is continuous for $d_{\fsymlinksp}$ at $(u\opt,
%	\vec{\sigma}\opt)$.
%\end{assumption}
%The continuity of $\nabla_\theta^2 \poploss(\theta, \vec{\sigma})$ 
%%$\nabla_\theta^2 \poploss(\theta, \vec{\sigma})$ at
%%$(u\opt, \vec{\sigma}\opt)$ 
%allows us to develop local asymptotic normality.
%To see we may expect the assumption to hold, we give reasonably simple
%conditions sufficient for it in Lemma~\ref{lem:d-continuity-guarantees}, Appendix~\ref{sec:semi-supp}.
Then
the next master result for semi-parametric
approaches characterizes the asymptotic behavior of the
semi-parametric estimator $\spe$ using the variance function
\begin{align*}
  C_{m, \vec{\sigma}\opt} := \frac{1}{m} \cdot
  \frac{\frac{1}{m} \sum_{j = 1}^m \E[\sigma_j\opt(Z) (1 - \sigma_j\opt(Z))]}{
    (\frac{1}{m} \sum_{j = 1}^m
    \E[{\sigma_j\opt}'(Z)])^2}.
\end{align*} 
\begin{theorem}
  \label{theorem:semiparametric-master-informal}
  Let Assumption~\ref{assumption:x-decomposition} hold and assume $|Z| > 0$
  has nonzero and continuous density $p(z)$ on $(0, \infty)$.  Suppose the
  sequence $\{\vec{\sigma}_n\}$ is consistent, i.e.,
  \begin{equation*}
    \LtwoPbold{\vec{\sigma}_n(-Y\<X, u\opt\>)
      - \vec{\sigma}\opt(-Y\<X, u\opt\>)}
    \cp 0,
  \end{equation*}
  the mapping $(\theta, \vec{\sigma}) \mapsto
  \nabla_\theta^2 \poploss(\theta,
  \vec{\sigma})$ is continuous for $d_{\flink}$ at $(u\opt,
  \vec{\sigma}\opt)$, and $\Ep [\ltwo{X}^4] < \infty$.  Then $\sqrt{n}(\spe - u\opt)$ is
  asymptotically normal, and the estimator $\what{u}\sem_{n,m} =
  \spe / \ltwos{\spe}$ satisfies
  \begin{equation*}
    \sqrt{n} \prn{\what{u}\sem_{n,m} - u\opt}
    \cd \normal\prn{0, 	C_{m, \vec{\sigma}\opt}
      \prn{\proj_{u\opt}^\perp \Sigma \proj_{u\opt}^\perp}^\dagger}.
  \end{equation*}
\end{theorem}
\noindent
Notably, Theorem~\ref{theorem:semiparametric-master-informal} exhibits optimal
$1/m$ covariance scaling whenever $\E[{\sigma\opt_j}'(Z)] \gtrsim 1$.

We give two example applications of the general
theory in Appendix~\ref{sec:semi-supp}: the first analyzing a full pipeline
for a single index model setting, which robustly estimates direction $u\opt$,
the link $\sigma\opt$, and then re-estimates $\theta\opt$; the second
assuming a stylized black-box crowdsourcing mechanism that provides
estimates of labeler reliability, highlighting how even in some
crowdsourcing scenarios, there could be substantial advantages to using full
label information.

\paragraph{Future work.}

In spite of the technical detail we require to prove our results,
we view this work as almost preliminary and hope that it inspires further
work on the full pipeline of statistical machine learning,
from dataset creation to model release.
%
Many questions remain.

On the theoretical side, we focus on a stylized model of
label aggregation, with majority vote---with the exception of the
crowdsourcing model above--functioning as the
stand-in for more sophisticated aggregation strategies. It seems challenging
to show that \emph{no} aggregation strategy can work as well as multi-label
strategies; it would be interesting to delineate the benefits
and drawbacks of more sophisticated denoising.
% We
% focus throughout on low-dimensional asymptotics, using asymptotic normality
% to compare estimators. 
Additionally, the insights make low-dimensional asymptotic predictions;
investigating high-dimensional scaling might yield new perspectives
as, for, classification datasets often become
separable~\cite{CandesSu19pnas}.
%%, which is consistent with modern
%%applied machine learning~\cite{ZhangBeHaReVi21}.
%but makes the asymptotics we
% derive impossible. One option
% is to investigate the asymptotics of maximum-margin estimators,
% which coincide with the limits of ridge-regularized logistic
% regression~\cite{RossetZhHa04, SoudryHoNaGuSr18}.

On the methodological and applied side, given the extent to which the
test/challenge dataset methodology drives progress in machine
learning~\cite{Donoho17}, it seems that developing newer datasets to
incorporate labeler uncertainty could yield substantial benefits.  A
particular refrain is that modern deep learning methods are overconfident in
their predictions~\cite{GoodfellowShSz15}; perhaps by
calibrating them to \emph{labeler} uncertainty we could
improve their robustness and performance.
%
Large-scale
models now frequently use only distantly supervised
labels~\cite{RadfordKiHaRaGoAgSaAsMiClKrSu21}, making
the study of appropriate data cleaning and collection more timely.
We look forward to deeper
investigations of the foundations of the full
practice of statistical machine learning.
